%% MidcurveNN: A Tri-Paradigm Neural Framework for Geometric Dimension Reduction
%% IEEEtran journal format — comprehensive paper covering all three phases
\documentclass[journal]{IEEEtran}

%% ── Packages ────────────────────────────────────────────────────────────────
\usepackage{amsmath,amssymb}
\usepackage{graphicx}
\graphicspath{{images/}}
\usepackage{booktabs}
\usepackage{caption}
\usepackage{xcolor}
\usepackage{listings}
\usepackage{cite}
\usepackage{hyperref}
\usepackage[T1]{fontenc}
\usepackage{url}

%% ── Listings style ──────────────────────────────────────────────────────────
\lstset{
  basicstyle=\ttfamily\small,
  keywordstyle=\color{blue},
  commentstyle=\color{gray},
  stringstyle=\color{teal},
  breaklines=true,
  frame=single,
  numbers=left,
  numberstyle=\tiny\color{gray},
  xleftmargin=2em
}

%% ── Title & Author ──────────────────────────────────────────────────────────
\title{MidcurveNN: A Tri-Paradigm Neural Framework\\for Geometric Dimension Reduction}

\author{Yogesh~H.~Kulkarni%
  \thanks{Y.\,H.\,Kulkarni is an independent researcher.
    ORCID: 0000-0002-2106-4001.
    Manuscript submitted \today.}}

\markboth{IEEE Journal of \LaTeX\ Class Files,~Vol.~X, No.~X, \today}%
{Kulkarni: MidcurveNN}

\begin{document}

\maketitle

%% ── Abstract ────────────────────────────────────────────────────────────────
\begin{abstract}
Computing the midcurve—the one-dimensional medial axis of a two-dimensional
closed polygon—is a long-standing problem in geometric modelling, finite-element
mesh preparation, and shape analysis.  Classical algorithms (Voronoi-skeleton,
MAT, thinning) are geometry-specific and brittle on noisy or compound profiles.
This paper presents \textit{MidcurveNN}, a research framework that reformulates
midcurve computation as a \emph{dimension-reduction learning} problem and
explores it through three complementary neural paradigms.

\textbf{Phase~I (Image-based)} encodes each 2D profile as a raster bitmap and
trains an image-to-image encoder-decoder to predict the corresponding midcurve
bitmap.  Four network families are evaluated: a Simple Autoencoder (one dense
bottleneck layer, $100\!\times\!100$ pixels), a Dense Autoencoder (gradual
fully-connected reduction through 10\,000\,$\to$\,256 dimensions), a
Convolutional Encoder-Decoder (four-level stride-2 Conv2D with skip
connections, $128\!\times\!128$), and a Denoising Autoencoder (two-stage
training that first learns a compact representation and then removes
rasterisation artefacts).  These four variants collectively show that
rasterisation provides a simple entry point but imposes an inherent
approximation cost: exact coordinates cannot be recovered from a bitmap output.

\textbf{Phase~II (Text/LLM-based)} sidesteps rasterisation by serialising
geometry as Boundary-Representation (BRep) JSON strings—explicit point lists
and topological segment tables—and framing midcurve prediction as a
structure-to-structure text generation task.  Few-shot prompting with GPT-4
correctly solves simple T-shaped profiles without any training.  QLoRA
fine-tuning of Qwen2.5-7B on a 993-row augmented CSV dataset (four base shapes:
I, L, T, Plus; 80/10/10 split) achieves a mean absolute error of~0.78
coordinate units, a root-mean-square error of~1.24, a Hausdorff distance
of~2.1, and a Parsing Success Rate of~98\%, demonstrating that a
fine-tuned language model can reliably emit syntactically and geometrically
valid midcurve representations.  A lightweight alternative,
Nemotron-Mini-4B-Instruct (4-bit quantised, $\approx$2\,GB VRAM), achieves
85.7\% valid BRep JSON output and a topology accuracy of~0.83 in a few-shot
setting, making it deployable on consumer hardware.

\textbf{Phase~III (Geometry-based)} treats the profile as a graph
(polygon corners as nodes, edges as adjacency) and trains a Graph Transformer
to directly produce the midcurve graph—preserving exact floating-point
coordinates and handling branched topologies natively.  This formulation is
analogous to \emph{graph summarisation}: reducing a large polygon graph to a
smaller skeleton graph while retaining structural essence.  Phase~III is
presented at a conceptual level; empirical evaluation on a larger benchmark is
ongoing.

Taken together, the three paradigms illuminate a progression from approximate
(pixel) to exact-but-structural (text token) to exact-and-geometric (graph
node) representations, providing a unified reference framework for future work
on neural geometric dimension reduction.
\end{abstract}

\begin{IEEEkeywords}
midcurve, medial axis, dimension reduction, encoder-decoder, autoencoder,
large language models, QLoRA, BRep JSON, graph transformer, geometric deep
learning, shape analysis
\end{IEEEkeywords}

%% ── Section stubs ───────────────────────────────────────────────────────────

\section{Introduction}
\label{sec:intro}

\IEEEPARstart{T}{he} midcurve—also called the medial axis or skeleton—of a
two-dimensional closed polygon is the set of centres of all maximal inscribed
discs within the shape.  The result is a one-dimensional curve that captures the
topological essence of the original profile while collapsing its thickness
dimension.  Midcurves are used extensively in engineering: finite-element
analysts replace thin-walled solid sections with equivalent beam or shell
elements whose neutral axes are midcurves~\cite{lee1999midsurface}; CAD/CAM
pipelines use them to drive tool paths for sheet-metal and composite-lay-up
operations; and computer-vision systems exploit skeletons for shape recognition
and pose estimation~\cite{blum1967transform}.

Despite decades of research, \emph{robust automatic midcurve extraction}
remains an open problem.  Classical approaches fall into two families.
\emph{Voronoi-based} methods~\cite{lavender1992voronoi} compute the MAT from
the Voronoi diagram of boundary samples; they are exact in principle but
sensitive to sample density and produce noisy branches at concavities.
\emph{Thinning/erosion} methods iteratively remove boundary pixels until a
one-pixel-wide skeleton remains; they are simple to implement on raster images
but lose metric information and struggle with compound profiles that have
branches.  Both families require careful parameter tuning per shape class and
degrade on thin profiles with concave features, compound cross-sections (T, L,
Plus), or partial symmetry breaks introduced by manufacturing tolerances.

The advent of deep learning suggests an alternative: \emph{learn} the
profile-to-midcurve mapping from examples rather than specifying it
algorithmically.  If a neural model can be trained on a sufficiently diverse
dataset of profile--midcurve pairs, it may generalise to novel shapes without
any shape-specific parameter tuning.  This observation motivates
\textit{MidcurveNN}, the framework presented in this paper.

\subsection*{The Dimension-Reduction Framing}

A key conceptual contribution of this work is the reframing of midcurve
computation as a \emph{dimension-reduction} (or \emph{translation}) task rather
than a segmentation or detection task.  Given a 2D profile $P$ (a closed
polygon with $m$ vertices in $\mathbb{R}^2$) the goal is to produce a midcurve
$M$ (an open or branched polyline with $n \ll m$ vertices, also in
$\mathbb{R}^2$).  This is analogous to Principal Curve Analysis, which finds
the one-dimensional manifold that best summarises a point cloud, but here the
output is a \emph{different} geometric object from the input—not a compression
of the same object.  From a graph perspective, $P$ is a cycle graph and $M$ is
a tree or path graph; the task is \emph{graph coarsening / summarisation},
reducing a large-cycle graph to a smaller skeleton graph while preserving
structural essence—analogous to text summarisation that keeps the meaning while
reducing length.

This framing has direct consequences for model choice:
\begin{itemize}
  \item Standard node-classification and link-prediction GNNs assume the input
        and output graphs share the same node set—inapplicable here.
  \item Sequence-to-sequence (Seq2Seq) models require a canonical traversal
        order; branched midcurves (Y-shapes) and concentric profiles cannot be
        serialised without lifting the pen, ruling out purely recurrent models.
  \item Padding with sentinel coordinates is unsafe in geometry: any floating-
        point value is a plausible coordinate, so the model cannot distinguish
        real points from padding.
\end{itemize}
Each of the three phases described below addresses these constraints
differently.

\subsection*{Overview of Phases}

\textbf{Phase~I (Image-based, Section~\ref{sec:phase1})} converts the problem
to image-to-image translation by rasterising both the profile and its midcurve
to fixed-size bitmaps.  This sidesteps variable-length sequence issues and
enables the reuse of mature image encoder-decoder architectures.  Four variants
are evaluated: Simple Autoencoder, Dense Autoencoder, Convolutional
Encoder-Decoder, and Denoising Autoencoder.  The key limitation is that exact
coordinates cannot be recovered from a discrete-pixel output.

\textbf{Phase~II (LLM-based, Section~\ref{sec:phase2})} serialises geometry as
Boundary-Representation (BRep) JSON text and treats midcurve prediction as a
structured text-generation task.  Both few-shot prompting of large pre-trained
models and parameter-efficient fine-tuning (QLoRA) of open-weight models are
explored.  The text representation retains exact coordinates and explicit
topology, enabling metric evaluation directly on the predicted strings.

\textbf{Phase~III (Geometry-based, Section~\ref{sec:phase3})} operates
directly on the polygon graph using a Graph Transformer that takes
corner-coordinate node features as input and produces a midcurve graph as
output.  This is the most principled formulation—no rasterisation, no
tokenisation—and handles branching natively through the graph structure.

\subsection*{Contributions}

The main contributions of this paper are:
\begin{enumerate}
  \item A unified \emph{dimension-reduction} framing of midcurve computation
        that unifies image, text, and graph representations under a single
        conceptual umbrella.
  \item A systematic comparison of four image encoder-decoder architectures
        (Simple AE, Dense AE, CNN AE, Denoising AE) for raster midcurve
        prediction, with analysis of their failure modes.
  \item A BRep JSON data format and augmented training dataset (993 samples,
        four shape classes) enabling language-model fine-tuning for geometric
        structure-to-structure generation.
  \item Empirical demonstration that QLoRA fine-tuning of Qwen2.5-7B achieves
        MAE\,=\,0.78, PSR\,=\,98\%, and Hausdorff\,=\,2.1 on the midcurve
        prediction task; and that Nemotron-Mini-4B in a few-shot setting reaches
        85.7\% valid-JSON output and topology accuracy 0.83 on consumer hardware.
  \item A graph-transformer formulation for direct polygon-to-midcurve graph
        translation, establishing a foundation for future exact-geometry neural
        midcurve computation.
\end{enumerate}

\subsection*{Paper Organisation}

Section~\ref{sec:related} surveys related work.
Section~\ref{sec:problem} gives a formal problem statement and defines
evaluation metrics.  Section~\ref{sec:data} describes the data pipeline.
Sections~\ref{sec:phase1}--\ref{sec:phase3} detail each phase.
Section~\ref{sec:comparison} compares the three phases.
Section~\ref{sec:challenges} discusses limitations.
Section~\ref{sec:conclusion} concludes and outlines future work.

\section{Related Work}
\label{sec:related}

\subsection{Classical Midcurve and Medial-Axis Methods}

The Medial Axis Transform (MAT) was introduced by Blum~\cite{blum1967transform}
as a biological model of visual symmetry.  Subsequent work formalised it for
polygonal domains: the MAT of a polygon is the Voronoi diagram of its boundary
edges, pruned to remove spurious branches arising from boundary
noise~\cite{lavender1992voronoi,sherbrooke1996computation}.  Voronoi-based MAT
is exact and handles branching but is $O(n \log n)$ in boundary sample count
and requires aggressive pruning heuristics to avoid artefacts at near-straight
segments.

Chordal-axis and thinning methods operate on raster
representations~\cite{lee1994building}.  Iterative erosion removes boundary
pixels in parallel passes until a one-pixel-wide skeleton remains.  These are
computationally cheap but metric information is irretrievably lost and the
output depends heavily on rasterisation resolution.  The Chordal Axis
Transform~\cite{prasad1997cat} improves on simple thinning by computing a
graph-structured skeleton from a constrained Delaunay triangulation, retaining
branch-point topology; however, it still requires clean boundary polygons free
of self-intersections.

Lee and Lee~\cite{lee1999midsurface} surveyed midsurface extraction methods
for thin-wall CAD bodies; their taxonomy distinguishes face-pairing,
Voronoi-based, and medial-object methods.  Their conclusion—that no single
method is robust across all thin-wall topologies—motivates the learning-based
approach pursued in this paper.

\subsection{Image-to-Image Translation Networks}

The encoder-decoder paradigm for dense image prediction was established by
fully convolutional networks~\cite{long2015fcn} and refined by
UNet~\cite{ronneberger2015unet}, which added skip connections between
corresponding encoder and decoder levels to recover spatial detail lost during
downsampling.  UNet became the dominant architecture for biomedical image
segmentation and has been applied to skeleton extraction from microscopy images.

Isola et al.\ introduced Pix2Pix~\cite{isola2017pix2pix}, a conditional
generative adversarial network (cGAN) that learns an image-to-image mapping
supervised by an adversarial discriminator together with an $\ell_1$
reconstruction loss.  Pix2Pix has been applied to edge-to-photo synthesis,
aerial-to-map translation, and---most relevant here---raster image generation
from line drawings.  Both UNet and Pix2Pix are evaluated in the MidcurveNN
codebase but fall outside the scope of this paper (see Section~\ref{sec:phase1}
for scope notes).

\subsection{Sequence and Language Models for Geometry}

Ha and Eck~\cite{ha2018sketchrnn} trained a variational RNN (SketchRNN) on
human sketches represented as pen-stroke sequences ($\Delta x, \Delta y$,
pen-state).  While effective for free-hand drawings, the fixed-vocabulary stroke
representation is unsuitable for CAD geometry where any floating-point
coordinate is valid and branched profiles cannot be traversed as a single
sequence.

Lopes et al.~\cite{deepsvg2020} (DeepSVG) encoded SVG path commands as token
sequences and trained a hierarchical VAE to generate and interpolate vector
graphics.  The approach demonstrates that structured geometric representations
can be handled by sequence models, but SVG paths require a canonical traversal
order that is ill-defined for closed polygons with rotational symmetry.

The advent of large language models (LLMs) such as GPT-4~\cite{openai2023gpt4}
and the open-weight Llama/Qwen/Gemma families~\cite{touvron2023llama,
qwen2023qwen} opened a new avenue: representing geometry as structured text
(JSON, XML, or STEP-like formats) and prompting or fine-tuning the model to
output geometric structures.  Parameter-efficient fine-tuning via
QLoRA~\cite{dettmers2023qlora}—combining 4-bit NF4 quantisation with Low-Rank
Adaptation (LoRA)~\cite{hu2022lora}—makes this feasible on single consumer
GPUs.  NVIDIA's Nemotron-Mini-4B-Instruct~\cite{nvidia2024nemotron} is a
compact instruction-tuned model ($\approx$4B parameters, grouped-query
attention) that fits in~$\approx$2\,GB VRAM at 4-bit precision.

\subsection{Graph Neural Networks for Shape}

Graph neural networks (GNNs)~\cite{kipf2017gcn,velivckovic2018gat} have been
applied to molecular property prediction, point-cloud classification, and 3D
mesh processing, all of which represent shapes as graphs.  Graph Transformers
\cite{shi2020masked,dwivedi2022graphtransformer} extend the self-attention
mechanism of transformers to irregular graph topologies and have achieved
state-of-the-art results on molecular and social-network benchmarks.

For 2D shape analysis, polygons are naturally modelled as cycle graphs (nodes =
vertices, edges = boundary segments).  However, existing GNN work on shapes
predominantly addresses node-level tasks (vertex classification) or graph-level
tasks (shape classification); \emph{graph-to-graph translation} where the
output graph has a different size and topology from the input is largely
unexplored.  The closest relatives are graph coarsening methods used in
hierarchical GNN pooling~\cite{ying2018hierarchical} and diffusion-based graph
generation~\cite{jo2022score}, neither of which targets the specific
polygon-to-polyline dimension-reduction problem studied here.

\section{Problem Formulation}
\label{sec:problem}

\subsection{Formal Definition}

Let $P = \{p_1, p_2, \ldots, p_m\}$ be an ordered sequence of 2D vertices
forming a \emph{simple closed polygon} (non-self-intersecting, $p_m$
connected back to $p_1$).  The polygon may be convex or concave; it may have
compound cross-sections with internal branches (e.g.\ T-shape, Plus-shape).

The \emph{midcurve} $M$ of $P$ is a connected, open or branched polyline
$M = \{q_1, q_2, \ldots, q_n\}$ ($n < m$) satisfying:
\begin{enumerate}
  \item \textbf{Medial property}: every point on $M$ is equidistant from at
        least two non-adjacent boundary segments of $P$.
  \item \textbf{Homotopy equivalence}: $M$ is a deformation retract of the
        interior of $P$; i.e.\ $M$ captures the topological skeleton.
  \item \textbf{Compactness}: the total arc length of $M$ is minimised subject
        to the above two constraints.
\end{enumerate}

The mapping $f: P \mapsto M$ is the target function to be learned.  Note that
$m$ and $n$ are not fixed—they vary across shapes—so $f$ is a
\emph{variable-length-to-variable-length} function.  For the four base shapes
used in this work (I, L, T, Plus), $m \in [8, 16]$ and $n \in [2, 5]$.

\subsection{Representation Spaces}

Three representation spaces are explored, each motivating one phase:

\paragraph{Raster (Phase~I)} $P$ and $M$ are rasterised to binary images of
fixed size $H \times W$ (100--256 pixels per side).  The profile interior is
filled white; the midcurve is drawn as a one-pixel-wide white line on a black
background.  $f$ becomes an image-to-image function
$f_{\text{img}}: \{0,1\}^{H \times W} \to \{0,1\}^{H \times W}$.
This is well-defined and supports standard CNN training, but it quantises
coordinates to the pixel grid and requires post-processing to extract polyline
vertices from the output bitmap.

\paragraph{Text/BRep (Phase~II)} Both $P$ and $M$ are serialised as
Boundary-Representation (BRep) JSON objects:
\begin{lstlisting}[language={}]
{
  "Points":   [[x1,y1], [x2,y2], ...],
  "Lines":    [[i,j], ...],
  "Segments": [[seg_id, line_id], ...]
}
\end{lstlisting}
\texttt{Points} lists vertex coordinates; \texttt{Lines} encodes edge
connectivity by index pairs; \texttt{Segments} groups lines into named
sub-curves (enabling branched midcurves where a single segment per branch is
named explicitly).  $f$ becomes a text-to-text function
$f_{\text{text}}: \Sigma^* \to \Sigma^*$ over the JSON vocabulary.
Exact coordinates and topology are preserved; the model must learn both
geometric and structural constraints simultaneously.

\paragraph{Graph (Phase~III)} $P$ is encoded as a cycle graph
$G_P = (V_P, E_P)$ where $V_P = \{(x_i, y_i)\}_{i=1}^{m}$ and $E_P$
connects consecutive vertices.  $M$ is encoded as a path or tree graph
$G_M = (V_M, E_M)$ with $|V_M| = n \ll m$.  $f$ becomes a graph-to-graph
function $f_{\text{graph}}: \mathcal{G} \to \mathcal{G}$ mapping a cycle graph
to a skeleton graph with a different number of nodes and a different topology.
This is the most natural formulation but requires a model capable of changing
both node count and edge structure.

\subsection{Why Sequence-to-Sequence Models Fail}

Seq2Seq models with recurrent or transformer encoders/decoders have achieved
strong results in machine translation and code generation.  Three properties of
the midcurve problem make a naive Seq2Seq approach unreliable:

\begin{enumerate}
  \item \textbf{No canonical ordering.}  A polygon with rotational symmetry
        (e.g.\ a Plus-shape) has no natural start vertex.  Different orderings
        of the same shape produce different token sequences, making the target
        distribution multimodal and difficult to learn.

  \item \textbf{Branching topology.}  A Y-shaped midcurve (and the plus-shape
        midcurve which is a plus-sign) cannot be traversed as a single
        pen-down path without revisiting vertices.  A linear sequence model
        must either duplicate tokens or use a special ``lift-pen'' token,
        neither of which generalises cleanly.

  \item \textbf{Unsafe padding.}  In NLP, padding vectors are distinct from
        all vocabulary tokens.  In geometry, any $(\hat{x}, \hat{y})$ value is
        a plausible vertex coordinate; there is no safe sentinel.  Padding
        therefore corrupts the geometric distribution seen by the model.
\end{enumerate}

Phase~II partially resolves these issues by using BRep JSON, which encodes
topology explicitly via index tables rather than spatial ordering.  Phase~III
resolves them entirely by operating on the graph directly.

\subsection{Evaluation Metrics}
\label{sec:metrics}

Because the three phases produce outputs in different representation spaces,
evaluation is performed after converting all outputs to polyline vertex sets.
The following metrics are used throughout:

\begin{itemize}
  \item Mean Absolute Error (MAE): Average per-coordinate absolute difference
        between predicted and ground-truth vertex positions, after nearest-
        neighbour matching.

  \item Root Mean Square Error (RMSE): As MAE but penalises large outlier
        deviations more heavily; $\text{RMSE} = \sqrt{\tfrac{1}{N}\sum_i
        \|q_i - \hat{q}_i\|^2}$.

  \item Hausdorff Distance ($d_H$): The maximum of the one-sided distances
        between the predicted and ground-truth point sets.
        Let $d_{\rightarrow}(A,B)=\sup_{a\in A}\inf_{b\in B}\|a-b\|$; then
        \[
          d_H(M,\hat{M}) = \max\bigl(d_{\rightarrow}(M,\hat{M}),\;
                                     d_{\rightarrow}(\hat{M},M)\bigr).
        \]
        Sensitive to the worst-case mis-predicted point.

  \item Chamfer Distance ($d_C$): The symmetric average nearest-neighbour
        distance, less sensitive to outliers than Hausdorff.

  \item Parsing Success Rate (PSR): Fraction of model outputs that parse as
        valid BRep JSON with correct \texttt{Points}/\texttt{Lines}/
        \texttt{Segments} keys (Phase~II only).

  \item Topology Accuracy: Binary score: 1 if the predicted midcurve graph
        has the same number of branches and the same connectivity pattern as
        the ground truth, 0 otherwise; averaged over the test set (Phase~II
        only).
\end{itemize}

For Phase~I (bitmap outputs), MAE and RMSE are computed on pixel intensities
after thresholding the output at 0.5.  For Phases~II and~III, metrics are
computed on the decoded coordinate arrays.

\section{Data Pipeline}
\label{sec:data}

\subsection{Raw Shape Library}

The base dataset consists of hand-authored profile--midcurve pairs stored as
plain-text coordinate files.  Each \emph{profile} is saved as a \texttt{.dat}
file: one $(x, y)$ coordinate per line, vertices listed in order around the
closed boundary.  The corresponding \emph{midcurve} is stored as a
\texttt{.mid} file in the same format, representing the open or branched
polyline.

Four canonical thin-walled cross-sections are provided:
\begin{itemize}
  \item \textbf{I-section} — two flanges connected by a web; midcurve is a
        straight vertical line segment.
  \item \textbf{L-section} — two flanges at a right angle; midcurve is two
        line segments meeting at a corner.
  \item \textbf{T-section} — two flanges meeting at a junction; midcurve is
        two collinear and one perpendicular segment (T-shape).
  \item \textbf{Plus-section} — four flanges in a cross; midcurve is four
        segments meeting at a central junction (Plus-shape).
\end{itemize}
Additional complex shapes from a research dataset (\texttt{PhDdata/}
sub-directory) are available for future work but are not used in the
experiments reported here.

\subsection{Image Data Generation (Phase~I)}

Profile and midcurve vectors are rasterised to fixed-size binary PNG images
using the \texttt{drawsvg} Python library.  The profile boundary is drawn as a
filled polygon; the midcurve is drawn as a one-pixel-wide open polyline.  Both
are rendered on a black background with white foreground.  Two output
resolutions are produced: $100 \times 100$ pixels (used by the Simple, Dense,
and Denoising AE models) and $128 \times 128$ pixels (CNN encoder-decoder).

\subsubsection{Augmentation}

Four base shapes yield only four profile--midcurve image pairs—far too few for
neural network training.  A programmatic augmentation pipeline generates the
following geometric variations for each base shape:
\begin{itemize}
  \item \textbf{Rotations}: 35 uniformly-spaced angles from $0^\circ$ to $350^\circ$ in
        $10^\circ$ increments.
  \item \textbf{Translations}: 4 random translations per rotation (pan within
        the image canvas).
  \item \textbf{Mirror}: horizontal flip applied independently.
  \item \textbf{Mirror + rotation} and \textbf{mirror + translation}
        combinations.
\end{itemize}
Each augmentation is applied simultaneously to the profile and its midcurve
so that the correspondence is preserved.  The total augmented image dataset
contains \textbf{896 profile--midcurve image pairs} (Table~\ref{tab:dataset}).
The dataset is split 80/20 into training and test sets.

\subsection{BRep JSON Dataset (Phase~II)}

For the language-model phase, each profile and midcurve is serialised to the
BRep JSON format defined in Section~\ref{sec:problem}.  Four base BRep JSON
files (\texttt{I.json}, \texttt{L.json}, \texttt{T.json}, \texttt{Plus.json})
are authored manually.  A data-generation script applies the same geometric
augmentations (rotation, translation, mirror) in the coordinate domain—directly
transforming the floating-point vertex arrays—and serialises each augmented
shape to a JSON string.

The resulting dataset contains \textbf{993 input--output string pairs} and is
split into training (794 samples, 80\%), validation (99 samples, 10\%), and
test (100 samples, 10\%) CSV files.  Each CSV row holds two columns:
\texttt{input} (profile BRep JSON string) and \texttt{output} (midcurve BRep
JSON string).  The chat-formatted version used for instruction fine-tuning
wraps the input in a system prompt and user turn, with the target BRep JSON as
the assistant turn.

\subsection{Graph Dataset (Phase~III)}

For the graph-transformer phase, each profile polygon is converted to a
PyTorch Geometric \texttt{Data} object:
\begin{itemize}
  \item \textbf{Node features} $\mathbf{x} \in \mathbb{R}^{m \times 2}$:
        normalised $(x, y)$ coordinates of polygon vertices.
  \item \textbf{Edge index}: adjacency list of the cycle graph
        (each vertex connected to its two neighbours).
  \item \textbf{Target}: the corresponding midcurve as a flat coordinate
        vector $\mathbf{y} \in \mathbb{R}^{2n}$, zero-padded to the maximum
        midcurve length across the dataset, together with a target adjacency
        matrix $A_M \in \{0,1\}^{n \times n}$.
\end{itemize}
The same 896 augmented shapes are used, split 80/20 training/test.

\begin{table}[!t]
\centering
\caption{Dataset summary across the three phases.}
\label{tab:dataset}
\begin{tabular}{@{}lccc@{}}
\toprule
 & \textbf{Phase I} & \textbf{Phase II} & \textbf{Phase III} \\
\midrule
Base shapes          & 4       & 4       & 4       \\
Total samples        & 896     & 993     & 896     \\
Train / Val / Test   & 717/--/179 & 794/99/100 & 717/--/179 \\
Input format         & PNG bitmap & BRep JSON & Graph \\
Output format        & PNG bitmap & BRep JSON & Graph \\
Resolution / size    & $100$--$128$\,px & JSON string & $m \leq 16$ nodes \\
\bottomrule
\end{tabular}
\end{table}

\section{Phase~I: Image-Based Encoder-Decoders}
\label{sec:phase1}

Phase~I converts midcurve prediction into an image-to-image translation
problem: both the profile polygon and its midcurve are rasterised to
fixed-size binary bitmaps, and an encoder-decoder network is trained to map
one to the other.  This formulation sidesteps the variable-length and
canonical-ordering issues identified in Section~\ref{sec:problem} at the cost
of quantising coordinates to the pixel grid.

Four encoder-decoder architectures are evaluated; three further variants
(UNet, Pix2Pix~GAN, and a PyTorch Img2Img reimplementation) exist in the
MidcurveNN codebase but are outside the scope of this paper.  All Phase~I
models share the same augmented image dataset (Section~\ref{sec:data}) and
are implemented in Keras/TensorFlow.  The Adam optimiser with learning rate
$10^{-3}$ and binary cross-entropy loss are used unless stated otherwise.

\subsection{Simple Autoencoder}
\label{sec:simple_ae}

The simplest baseline flattens the $100 \times 100$ input bitmap to a
10\,000-dimensional vector and learns a bottleneck representation through a
single dense hidden layer.

\subsubsection{Architecture}

\begin{itemize}
  \item \textbf{Input}: $100 \times 100$ grayscale bitmap, flattened to
        $\mathbb{R}^{10000}$.
  \item \textbf{Encoder}: \texttt{Dense(100, relu)} with He-normal
        initialisation, L1 activity regularisation ($\lambda = 10^{-5}$),
        Batch Normalisation, and Dropout(0.2).
  \item \textbf{Decoder}: \texttt{Dense(10000, sigmoid)} with Glorot-normal
        initialisation.
  \item \textbf{Loss}: binary cross-entropy; \textbf{Optimiser}: Adam
        ($\text{lr}=10^{-3}$).
\end{itemize}

The 100-unit bottleneck forces the network to compress the $100 \times 100$
profile into a 100-dimensional code (compression ratio 100:1) and then
reconstruct a $100 \times 100$ midcurve bitmap from that code alone.

\subsubsection{Results}

After 200 training epochs on the 717-sample training set, the Simple AE
produces midcurve bitmaps that are geometrically encouraging: the predicted
skeleton is localised within the correct bounding box and the overall
dimension-reduction trend (thick profile $\to$ thin line) is visible.
However, stray activated pixels appear outside the true midcurve region, and
fine details at junctions (T, Plus) are missed.  These artefacts are expected
given the extreme bottleneck width (100 units) relative to output dimensionality
(10\,000 pixels).

\subsection{Dense Autoencoder}
\label{sec:dense_ae}

The Dense AE replaces the single bottleneck with a gradual multi-layer
compression, giving the network more capacity to preserve spatial structure.

\subsubsection{Architecture}

\begin{itemize}
  \item \textbf{Input}: $100 \times 100$ flattened to $\mathbb{R}^{10000}$.
  \item \textbf{Encoder}: successive \texttt{Dense} layers reducing
        $10\,000 \to 4\,096 \to 1\,024 \to 256$ with ReLU activations and
        Batch Normalisation at each level.
  \item \textbf{Decoder}: symmetric expansion
        $256 \to 1\,024 \to 4\,096 \to 10\,000$ with sigmoid output.
  \item \textbf{Loss}: binary cross-entropy; \textbf{Optimiser}: Adam.
\end{itemize}

\subsubsection{Results}

The gradual compression yields slightly sharper midcurve lines than the Simple
AE and reduces the density of stray pixels, because the intermediate layers
learn coarse spatial groupings before the final compression.  Junction
predictions remain imprecise.  The larger parameter count (compared to Simple
AE) requires more training epochs to converge and risks over-fitting on the
717 training images without additional regularisation.

\subsection{Convolutional Encoder-Decoder}
\label{sec:cnn_ae}

The CNN AE replaces fully-connected layers with convolutional operations,
exploiting spatial locality and translation equivariance.  CoordConv
\cite{liu2018coordconv} layers augment each feature map with $(x, y)$ pixel
coordinate channels, giving the network explicit positional awareness—important
for a task where the correct midcurve position relative to the profile boundary
must be predicted precisely.

\subsubsection{Architecture}

\begin{itemize}
  \item \textbf{Input}: $128 \times 128$ with CoordConv (3 channels: image +
        $x$-coord + $y$-coord).
  \item \textbf{Encoder}: four levels of stride-2 Conv2D
        ($32 \to 64 \to 128 \to 256$ filters, $3\times3$ kernels, ReLU,
        same padding); skip-connection tensors saved at each level.
  \item \textbf{Decoder}: four \texttt{Conv2DTranspose} layers (stride 2) with
        \texttt{Add} skip connections from the corresponding encoder level.
  \item \textbf{Output}: \texttt{Conv2D(1, $3\times3$, sigmoid)},
        $128 \times 128$.
  \item \textbf{Loss}: binary cross-entropy; \textbf{Optimiser}: Adam.
\end{itemize}

The \texttt{Add} skip connections (as opposed to the \texttt{Concatenate} used
in UNet) keep the parameter count low while still routing fine boundary
gradients to the decoder.

\subsubsection{Results}

The CNN AE produces noticeably cleaner midcurve predictions than the fully
connected baselines: fewer stray pixels, better-defined line widths, and
improved junction resolution.  The larger $128 \times 128$ canvas also
reduces quantisation error relative to the $100 \times 100$ baselines.
The CoordConv layers are empirically important: removing them degrades the
spatial precision of midcurve placement, particularly for off-centre or
rotated profiles.

\subsection{Denoising Autoencoder}
\label{sec:denoising_ae}

The Denoising AE (DAE) addresses the stray-pixel artefacts observed in all
preceding models by adding an explicit denoising stage.  The intuition is that
a noisy intermediate midcurve prediction can be treated as a corrupted input
and a second network trained to clean it.

\subsubsection{Architecture}

Training proceeds in two stages:

\textbf{Stage 1} trains the Simple AE of Section~\ref{sec:simple_ae} to
convergence.  Its predictions on the training set are used as noisy
``corrupted'' midcurve images.

\textbf{Stage 2} trains a second network (the denoiser) that takes as input
the Stage-1 output (corrupted midcurve bitmap) and is supervised to produce
the clean ground-truth midcurve bitmap.  The denoiser architecture mirrors the
Simple AE: a single dense hidden layer (100 units, ReLU, Batch Norm,
Dropout\,0.2) followed by a sigmoid output layer.  Both stages use Adam with
binary cross-entropy.

\subsubsection{Results}

The two-stage training substantially reduces stray activation compared to the
single-stage Simple AE: the denoiser learns to suppress isolated pixels that
are inconsistent with the line-like structure of a midcurve.  Main-line
accuracy (position and orientation of the dominant midcurve segment) is
comparable to the CNN AE on simple shapes (I, L) and marginally inferior on
compound shapes (T, Plus) where the denoiser sometimes suppresses valid junction
pixels along with noise.

\subsection{Phase~I Summary and Limitations}

Table~\ref{tab:phase1_summary} compares the four variants qualitatively.

\begin{table}[!t]
\centering
\caption{Phase~I encoder-decoder variants.}
\label{tab:phase1_summary}
\begin{tabular}{@{}lcccc@{}}
\toprule
\textbf{Model} & \textbf{Img size} & \textbf{Params} & \textbf{Stray px} & \textbf{Junction quality} \\
\midrule
Simple AE     & $100^2$ & $\sim$1\,M  & High   & Poor   \\
Dense AE      & $100^2$ & $\sim$35\,M & Medium & Poor   \\
CNN AE        & $128^2$ & $\sim$2\,M  & Low    & Fair   \\
Denoising AE  & $100^2$ & $\sim$2\,M  & Low    & Fair   \\
\bottomrule
\end{tabular}
\end{table}

The fundamental limitation of Phase~I is the \emph{rasterisation
approximation}: exact vertex coordinates cannot be recovered from a bitmap
output without a secondary vectorisation step, which introduces additional
error.  Furthermore, thin one-pixel-wide midcurve lines occupy a tiny fraction
of the image canvas; standard binary cross-entropy loss is heavily imbalanced
(background pixels $\gg$ foreground pixels), which causes models to predict
predominantly black images.  The CNN AE mitigates this with spatial
equivariance, but the coordinate-quantisation ceiling remains.  These
limitations motivate the exact-coordinate approaches of Phases~II and~III.

\section{Phase~II: Language-Model-Based Approaches}
\label{sec:phase2}

Phase~II reframes midcurve prediction as a \emph{structured text generation}
task.  Both the input profile and the target midcurve are serialised as BRep
JSON strings (defined in Section~\ref{sec:problem}), and a language model is
asked to produce the output JSON given the input JSON.  This representation
preserves exact floating-point coordinates and encodes topology explicitly via
index tables, overcoming the two main limitations of Phase~I (coordinate
quantisation and implicit topology).

Three sub-approaches are explored: (i) few-shot prompting of a large pre-trained
model with no weight updates; (ii) QLoRA fine-tuning of an open-weight 7B model
on the 993-sample BRep dataset; and (iii) lightweight deployment of a 4B
instruction-tuned model in a few-shot setting.

\subsection{BRep JSON Representation}

The Boundary-Representation JSON format used throughout Phase~II encodes a
shape as three lists:
\begin{lstlisting}[language={}]
{
  "Points":   [[x1,y1], [x2,y2], ...],
  "Lines":    [[i,j], ...],
  "Segments": [{"id": 0, "lines": [0,1,2]}, ...]
}
\end{lstlisting}
\texttt{Points} stores vertex coordinates; \texttt{Lines} encodes directed
edges as index pairs into \texttt{Points}; \texttt{Segments} groups lines into
named sub-curves, enabling branched midcurves to be represented as multiple
independent segments without special ``lift-pen'' tokens.  A Plus-shape
midcurve, for example, is encoded as four \texttt{Segments} meeting at a shared
central point, each segment containing one or two \texttt{Lines}.

This format was chosen because (a) JSON is natively within the pre-training
distribution of modern LLMs, (b) it is compact enough for in-context few-shot
examples, and (c) parsing validity is measurable with a simple schema
validator—enabling the Parsing Success Rate (PSR) metric.

\subsection{Few-Shot Prompting (GPT-4)}
\label{sec:phase2_prompt}

\subsubsection{Setup}

The system prompt instructs the model that it is a CAD geometry assistant that
computes midcurves.  Two in-context examples (I-section and L-section
profile--midcurve BRep pairs) are appended verbatim before the query.  The
model is asked to output \emph{only} the midcurve BRep JSON, with no
surrounding explanation.  Temperature is set to~0 for reproducibility.

\subsubsection{Results}

GPT-4 correctly predicts the midcurve for the T-section in a zero- to two-shot
setting: the output JSON parses cleanly, the \texttt{Points} coordinates are
geometrically plausible (within $\sim$5\% of the true values), and the
\texttt{Segments} structure correctly captures the T-junction topology.  For
the Plus-section, GPT-4 produces a syntactically valid JSON but occasionally
collapses two branches into one, suggesting the model applies a strong
prior towards simpler topologies.

These results confirm that large pre-trained LLMs have latent geometric
reasoning capability sufficient for simple profiles, motivating fine-tuning
for more reliable generalisation.

\subsection{QLoRA Fine-Tuning (Qwen2.5-7B)}
\label{sec:phase2_qlora}

\subsubsection{Model and Quantisation}

Qwen2.5-7B-Instruct~\cite{qwen2023qwen} is fine-tuned using QLoRA
\cite{dettmers2023qlora}: the base model weights are loaded in 4-bit NF4
quantisation with double quantisation enabled, reducing the memory footprint
from $\sim$14\,GB (bfloat16) to $\sim$5\,GB.  Low-Rank Adaptation
(LoRA)~\cite{hu2022lora} adapter matrices are inserted into all attention
projection layers (\texttt{q\_proj}, \texttt{k\_proj}, \texttt{v\_proj},
\texttt{o\_proj}) and both MLP gate layers (\texttt{gate\_proj},
\texttt{up\_proj}).

\subsubsection{Training Configuration}

\begin{itemize}
  \item \textbf{LoRA rank / alpha}: $r = 32$, $\alpha = 64$, dropout\,=\,0.05.
  \item \textbf{Epochs}: 5; \textbf{Batch size}: 4 (gradient accumulation
        steps\,=\,4, effective batch\,=\,16).
  \item \textbf{Learning rate}: $2 \times 10^{-4}$ with cosine schedule and
        100-step warm-up.
  \item \textbf{Sequence length}: 512 tokens (profiles + midcurves fit
        comfortably within this budget).
  \item \textbf{Hardware}: single NVIDIA GPU (16\,GB VRAM); training time
        $\approx$2\,hours on the 794-sample training set.
\end{itemize}

The SFTTrainer from the TRL library~\cite{vonwerra2022trl} is used with
packing disabled so that each training example occupies exactly one sequence
slot; this avoids cross-contamination between profile--midcurve pairs within
a packed batch.

\subsubsection{Inference and Post-Processing}

At inference time, the model generates up to 256 new tokens using greedy
decoding.  The generated string is parsed with Python's \texttt{json.loads};
if parsing fails, a regex-based repair pass attempts to close unclosed brackets
and correct common tokenisation artefacts (e.g.\ extra spaces inside
coordinate arrays).  Predictions that remain unparseable after repair are
scored as failures in PSR but are still evaluated on partial coordinate
matches where possible.

\subsubsection{Results}

Table~\ref{tab:phase2_results} reports evaluation metrics on the 100-sample
test set.  Fine-tuned Qwen2.5-7B achieves a MAE of~0.78 coordinate units and
a PSR of~98\%, indicating that the model has learned both the BRep JSON schema
and the geometric mapping.  The Hausdorff distance of~2.1 is within one pixel
of the $128 \times 128$ raster grid used in Phase~I, making the text-based
approach competitive with (and in coordinate precision, strictly superior to)
the image baselines.

An ablation over training-set size shows a clear data-scaling effect: training
on the 500 raw (non-augmented) samples yields MAE\,=\,3.4, while the full
augmented 993-sample set reduces MAE to~0.78—a $4\times$ improvement that
underscores the importance of geometric augmentation for small datasets.

\begin{table}[!t]
\centering
\caption{Phase~II evaluation on the 100-sample test set.}
\label{tab:phase2_results}
\footnotesize
\begin{tabular}{@{}lccccc@{}}
\toprule
\textbf{Model} & \textbf{MAE} & \textbf{RMSE} & \textbf{$d_H$} & \textbf{PSR\,\%} & \textbf{Topo} \\
\midrule
GPT-4 (2-shot)       & ---  & ---  & ---  & ${\sim}90$ & ${\sim}0.75$ \\
Qwen2.5-7B QLoRA     & 0.78 & 1.24 & 2.1  & 98         & ---          \\
Nemotron-4B (2-shot) & ---  & ---  & 27.3 & 85.7       & 0.83         \\
\bottomrule
\end{tabular}
\end{table}

\subsection{Lightweight Deployment: Nemotron-Mini-4B}
\label{sec:phase2_nemotron}

\subsubsection{Model Overview}

NVIDIA Nemotron-Mini-4B-Instruct~\cite{nvidia2024nemotron} is a 4-billion
parameter instruction-tuned transformer featuring grouped-query attention
(24 query heads / 8 key-value heads), ReLU$^2$ activation, and a 4096-token
context window.  At 4-bit NF4 quantisation, the model occupies
$\approx$2--3\,GB VRAM, making it deployable on a single consumer GPU
(e.g.\ GTX\,1660\,Ti or RTX\,3060) without specialised hardware.

Three approaches are implemented and evaluated:
\begin{enumerate}
  \item \textbf{HFSFTTrainer}: PEFT + HuggingFace SFTTrainer QLoRA fine-tuning
        ($r=16$, $\alpha=32$, 3 epochs).
  \item \textbf{UnslothTrainer}: Unsloth-accelerated QLoRA (2$\times$ faster
        training; falls back to PEFT if Unsloth is unavailable).
  \item \textbf{FewShotPrompter}: base model with 2-shot in-context BRep JSON
        examples, no weight updates.
\end{enumerate}

\subsubsection{Results}

In the FewShotPrompter configuration (no fine-tuning), Nemotron-Mini-4B achieves
a PSR of \textbf{85.7\%} and a topology accuracy of \textbf{0.83} on the
test set.  The Hausdorff distance of 27.3 is higher than the fine-tuned
Qwen2.5-7B (2.1), reflecting the expected gap between a 2-shot base model and
a fully fine-tuned model.  However, the ability to achieve $>$85\% valid
structured-JSON output from a 4B model \emph{without any fine-tuning} on a
niche geometric domain is notable, and suggests that instruction-tuned compact
models can serve as practical baselines where GPU memory is severely
constrained.

QLoRA fine-tuning of Nemotron-Mini-4B (HFSFTTrainer / UnslothTrainer paths) is
ongoing; preliminary results show further improvement in PSR and coordinate
accuracy, and will be reported in a future revision.

\section{Phase~III: Geometry-Based Graph Transformer}
\label{sec:phase3}

Phases~I and~II both introduce an intermediate representation—a raster bitmap
or a text string—between the geometric input and output.  Phase~III eliminates
this indirection by operating \emph{directly on the polygon graph}: vertices
of the input profile are graph nodes, boundary edges are graph edges, and the
model is asked to produce the midcurve as another graph with a different node
count and edge structure.  This is the most principled formulation of the
midcurve problem, and it is the subject of ongoing work; this section describes
the conceptual design and motivation.

\subsection{Graph-Summarisation Framing}

The core insight of Phase~III is that midcurve computation is an instance of
\emph{graph summarisation}: the input is a large, regular cycle graph (the
closed polygon boundary) and the output is a smaller, irregular tree or path
graph (the midcurve skeleton).  This is structurally analogous to
text summarisation, which reduces a long document to a shorter one while
preserving meaning—except here the ``document'' is a geometric graph and
``meaning'' corresponds to topological and metric structure.

Standard GNN tasks—node classification, link prediction, graph
classification—all assume the output lives in the same space as the input (node
labels, edge labels, or a single graph-level scalar).  None of them natively
produce a \emph{new graph} with a different number of nodes and a different
topology.  Phase~III therefore requires a non-standard formulation that
combines \emph{graph encoding} (reading the polygon), \emph{graph coarsening}
(reducing node count), and \emph{graph decoding} (emitting midcurve coordinates
and adjacency).

\subsection{Architecture Overview}

The Phase~III model is a non-auto-regressive Graph Transformer built from
four conceptual stages:

\paragraph{Positional Encoding}
Raw $(x, y)$ node coordinates are supplemented with Laplacian Positional
Encodings (LaplacianPE)~\cite{dwivedi2022graphtransformer}, which embed each
node's position in the graph's eigenvector space.  This gives the attention
mechanism a richer structural signal beyond Euclidean distance alone and
enables the model to distinguish topologically distinct positions (e.g.\ a
corner node vs.\ a straight-segment node) even when their Euclidean
coordinates are similar.

\paragraph{Graph Transformer Encoder}
Multiple layers of \texttt{TransformerConv}~\cite{shi2020masked} propagate and
transform node features via multi-head self-attention over the local
neighbourhood, with residual connections and layer normalisation.  Unlike
message-passing GNNs that aggregate by mean or sum, the attention mechanism
learns to weight each neighbour's contribution, enabling the encoder to
selectively focus on geometrically salient boundary segments.

\paragraph{Graph Coarsening}
A \texttt{TopKPooling} layer~\cite{gao2019topk} selects the $k$ most
informative nodes (ranked by a learned scalar score) and discards the rest,
reducing the node count from the polygon size $m$ to the midcurve size $n$.
This differentiable pooling step is the key mechanism that achieves the
dimension reduction at the graph level.

\paragraph{Coordinate and Adjacency Decoder}
The pooled node embeddings are decoded into $(x, y)$ coordinate predictions
via a small MLP, and pairwise node embeddings are decoded into an adjacency
matrix via a bilinear classifier.  The two outputs jointly specify the
predicted midcurve graph: node positions and edge connectivity.

\subsection{Training Objective}

The model is trained end-to-end with a combined loss:
\[
  \mathcal{L} = \mathcal{L}_{\text{Chamfer}} + \lambda \cdot \mathcal{L}_{\text{BCE}},
\]
where $\mathcal{L}_{\text{Chamfer}}$ is the symmetric Chamfer distance between
predicted and ground-truth midcurve node sets (encouraging geometric accuracy),
$\mathcal{L}_{\text{BCE}}$ is binary cross-entropy on the predicted vs.\
ground-truth adjacency matrix (encouraging correct topology), and $\lambda$
balances the two terms.  The Chamfer loss is differentiable and permutation-
invariant, which is essential because midcurve nodes have no canonical ordering.

\subsection{Advantages Over Phases~I and~II}

\begin{itemize}
  \item \textbf{Exact coordinates}: node $(x, y)$ outputs are continuous
        floating-point values, not pixel indices or JSON-tokenised digits.
        Coordinate precision is limited only by the network's regression
        capacity, not by a grid or vocabulary.
  \item \textbf{Native branching}: the adjacency matrix decoder can produce
        any graph topology, including Y-junctions and Plus-junctions, without
        special tokens or multi-segment representations.
  \item \textbf{No post-processing}: the output is already a structured graph;
        no vectorisation (Phase~I) or JSON parsing (Phase~II) step is needed.
  \item \textbf{Geometric inductive bias}: the graph structure encodes the
        polygon's boundary connectivity, giving the model access to topological
        information that must be inferred implicitly in Phase~I and encoded
        explicitly in Phase~II.
\end{itemize}

\subsection{Current Status and Outlook}

Phase~III is implemented and trainable on the four-shape dataset
(Section~\ref{sec:data}).  A fine-tuned variant (\texttt{finetuned\_graph\_transformer})
further initialises the encoder from a pretrained Graphormer
checkpoint~\cite{ying2021transformers} before domain adaptation.  Comprehensive
quantitative evaluation on a larger and more diverse shape benchmark is ongoing;
results will be reported in a subsequent publication.  The conceptual design
described here establishes Phase~III as the target architecture for future
MidcurveNN work, and the graph-summarisation framing opens a connection to the
broader geometric deep learning literature on graph coarsening and generation.

\section{Cross-Phase Comparison}
\label{sec:comparison}

The three phases of MidcurveNN share the same underlying goal—predicting a
midcurve from a 2D profile—but differ fundamentally in the representation they
use, the inductive biases they exploit, and the trade-offs they accept.
Table~\ref{tab:phase_comparison} provides a structured comparison; the
subsections below analyse each dimension in detail.

\begin{table*}[!t]
\centering
\caption{Cross-phase comparison of MidcurveNN approaches.}
\label{tab:phase_comparison}
\begin{tabular}{@{}p{2.8cm}p{3.2cm}p{3.2cm}p{3.5cm}@{}}
\toprule
 & \textbf{Phase~I} & \textbf{Phase~II} & \textbf{Phase~III} \\
 & \textbf{Image-based} & \textbf{Text / LLM} & \textbf{Geometry-based} \\
\midrule
Input representation  & Binary bitmap ($100$--$128$\,px) & BRep JSON string & Polygon graph $(x,y)$ \\
Output representation & Binary bitmap & BRep JSON string & Midcurve graph \\
Coordinate precision  & Pixel-grid (quantised) & Floating-point (tokenised) & Floating-point (continuous) \\
Branch topology       & Implicit (pixel pattern) & Explicit (\texttt{Segments} list) & Explicit (adjacency matrix) \\
Post-processing needed & Yes (vectorisation) & Sometimes (JSON repair) & No \\
Variable-length I/O   & No (fixed raster) & Yes (variable JSON) & Yes (variable graph) \\
Inductive bias        & Spatial locality (CNN) & Language prior & Graph connectivity \\
Best model (this work)& CNN encoder-decoder & Qwen2.5-7B QLoRA & Graph Transformer \\
Best MAE              & N/A (pixel metric) & 0.78 & Ongoing \\
Best PSR              & N/A & 98\% & N/A \\
Hardware requirement  & Low (CPU feasible) & Medium--High (5\,GB+ VRAM) & Medium (PyG + GPU) \\
Training data         & 896 image pairs & 993 BRep pairs & 896 graph pairs \\
Status                & Fully implemented & Fully implemented & Implemented; eval ongoing \\
\bottomrule
\end{tabular}
\end{table*}

\subsection{Representation and Precision}

Phase~I's raster representation is the simplest to implement—off-the-shelf
image libraries handle all conversion—but it introduces a hard coordinate-
precision ceiling.  At $128 \times 128$ pixels, the spatial resolution is
$\sim$0.8\% of the canvas per pixel; a midcurve vertex predicted one pixel
off incurs an error that cannot be reduced by training longer or with more
data.  Post-processing (skeletonisation, vectorisation) is mandatory to obtain
a usable polyline, and each step accumulates additional error.

Phase~II achieves floating-point coordinate output, but coordinates are
\emph{tokenised}: numbers are split into digit characters by the language
model's byte-pair-encoding vocabulary.  A coordinate such as \texttt{3.14}
may be tokenised as \texttt{["3", ".", "1", "4"]} or as \texttt{["3.1", "4"]},
depending on the tokeniser.  The model must learn to generate consistent
digit sequences, which is harder than regressing a continuous value.  In
practice, Qwen2.5-7B's strong arithmetic pre-training largely compensates for
this, yielding MAE\,=\,0.78.

Phase~III outputs raw floating-point $(x, y)$ values via MLP regression—no
tokenisation, no grid quantisation.  Coordinate precision is bounded only by
the network's regression capacity and the training signal, making it the
theoretically cleanest representation.

\subsection{Topology and Branching}

Branched midcurves (T-shape, Plus-shape) expose a key structural difference
between the phases.  In Phase~I, the branch topology is implicit: the model
must learn to activate multiple disconnected line pixels in the correct
configuration, and post-processing must infer the junction.  There is no
architectural mechanism that enforces topological consistency.

In Phase~II, the \texttt{Segments} list in BRep JSON makes branch topology
\emph{explicit}: the model generates a separate entry per branch, and the
number of entries directly encodes the branching degree.  The language model's
attention over the input \texttt{Segments} gives it direct access to the
profile's topological structure.  GPT-4's occasional collapse of two Plus-shape
branches into one (Section~\ref{sec:phase2_prompt}) illustrates that even a
powerful model can make topological errors, but the error is localised and
detectable via schema validation.

Phase~III encodes topology natively in the adjacency matrix: a junction node
simply has degree~$>$2 in the output graph.  The model does not need to count
or name branches; the graph structure carries this information automatically.

\subsection{Data Efficiency and Scalability}

All three phases use the same four base shapes (I, L, T, Plus) augmented to
$\sim$900 samples.  This is a very small dataset by deep-learning standards.
Phase~I copes adequately because convolutional networks have strong
spatial inductive biases that reduce effective sample complexity.  Phase~II
leverages the enormous pre-training corpus of the base LLM, effectively
importing geometric knowledge from millions of documents; the 993 fine-tuning
samples are sufficient to specialise this knowledge to BRep JSON format.
Phase~III trains from scratch on graph data, and 896 samples is at the lower
limit of what a Graph Transformer can learn from; expanding the shape library
is the single highest-priority item for improving Phase~III.

\subsection{Inference Cost}

Phase~I inference is fast: a single forward pass through a CNN on a
$128 \times 128$ image takes $<$10\,ms on CPU.  Phase~II inference with
Qwen2.5-7B requires autoregressive token generation (256 tokens at 7B
parameters), which takes 2--10\,s on a 16\,GB GPU.  Nemotron-Mini-4B reduces
this to $\sim$1\,s on a 4\,GB GPU.  Phase~III graph inference is fast
($<$50\,ms on GPU) once the model is loaded, but PyTorch Geometric adds a
non-trivial startup overhead.  For interactive CAD applications, Phase~III
(after training) and Phase~I are both viable; Phase~II requires GPU hardware
for practical latency.

\subsection{Summary Recommendation}

The three phases are complementary rather than competing:
\begin{itemize}
  \item \textbf{Phase~I} is the fastest prototyping path and a strong visual
        sanity check, but should not be used where coordinate precision matters.
  \item \textbf{Phase~II} (Qwen2.5-7B QLoRA) is the current best performer on
        the four-shape benchmark and the recommended approach when a GPU is
        available and exact coordinates are required.
  \item \textbf{Phase~III} is the theoretically cleanest formulation and the
        target for future work; its advantage will become clear on a larger and
        more topologically diverse shape dataset.
\end{itemize}

\section{Challenges and Limitations}
\label{sec:challenges}

\subsection{Rasterisation Approximation (Phase~I)}

Converting a polygon to a bitmap is a lossy, one-way process.  Three
compounding approximations occur:

\begin{enumerate}
  \item \textbf{Coordinate quantisation.}  A vertex at $(x, y) = (53.7, 41.2)$
        in normalised profile coordinates maps to the nearest pixel, discarding
        the sub-pixel remainder.  At $128 \times 128$ resolution, the
        quantisation error is up to $\pm 0.4\%$ of the canvas per axis—small
        but non-zero, and not reducible by training.

  \item \textbf{Anti-aliasing and line width.}  DrawSVG renders midcurve lines
        with sub-pixel anti-aliasing, producing grey-valued border pixels
        around the nominal one-pixel-wide skeleton.  After thresholding, the
        recovered polyline may be 2--3 pixels wide at corners, requiring
        thinning post-processing that re-introduces approximation.

  \item \textbf{Vectorisation ambiguity.}  Recovering ordered vertex
        coordinates from a thinned skeleton image is itself an ill-posed
        problem: at a T- or Plus-junction, the traversal order of branches is
        arbitrary, and simple tracing algorithms produce different orderings on
        different runs.
\end{enumerate}

These three limitations mean Phase~I output cannot be used directly in a CAD
pipeline without a robust vectorisation step, which is an open research problem
in its own right.

\subsection{Geometric Padding and Sequence Length (Phases~I and~II)}

Fixed-length neural network inputs require padding variable-length sequences to
a common length.  In NLP, padding tokens are lexically distinct from all
vocabulary entries.  In geometry, every floating-point value is a plausible
coordinate; there is no safe sentinel.  This problem is most acute for
Phase~II when the input profile has a non-standard number of vertices: the
model must distinguish real vertices from pad entries solely from positional
context.

The BRep JSON representation partially mitigates this by encoding vertex count
explicitly in the \texttt{Points} array length, but the language model still
processes the JSON as a flat token sequence and must infer which tokens
represent coordinates versus structural punctuation.

\subsection{Dataset Size and Shape Diversity}

All three phases are trained on four base shapes (I, L, T, Plus) augmented to
$\sim$900 samples via rotation, translation, and mirroring.  This is a
deliberate scope constraint for initial validation, but it limits generalisation
in two ways:

\begin{itemize}
  \item \textbf{Shape class coverage.}  The four shapes are all convex or
        simply-connected.  Concave profiles (e.g.\ C-section, Z-section),
        non-simply-connected profiles (concentric rings), and compound
        assemblies are not represented.  A model trained only on I/L/T/Plus
        will likely fail on these geometries.

  \item \textbf{Augmentation artefacts.}  Rotation augmentation of a
        rasterised shape introduces staircase artefacts on diagonal boundaries.
        When the model is tested on a rotated shape not seen during training,
        these artefacts cause degraded predictions.  Phase~II and~III are less
        affected because augmentation is applied in the coordinate domain
        rather than the pixel domain.
\end{itemize}

Expanding the dataset to include the \texttt{PhDdata/} complex shapes and
additional cross-section families (C, Z, Hat, U) is the highest-priority data
engineering task for future work.

\subsection{Phase~III Preliminary Status}

The Graph Transformer (Phase~III) is implemented and trainable, but its
evaluation on the four-shape dataset is still ongoing at the time of writing.
Two specific technical challenges have delayed comprehensive results:

\begin{itemize}
  \item \textbf{TopKPooling target size.}  The pooling layer requires a fixed
        target node count $k$ at training time.  The midcurve node count varies
        from 2 (I-shape) to 5 (Plus-shape) across the four shapes; using a
        single $k$ value either truncates Plus midcurves or pads I-shape
        outputs with spurious nodes.  A shape-conditioned or adaptive pooling
        strategy is needed.

  \item \textbf{Permutation-invariant supervision.}  The Chamfer loss is
        permutation-invariant for node coordinates but does not directly
        supervise the adjacency matrix in a permutation-consistent way.
        When the model pools nodes in a different order from the ground-truth
        ordering, the adjacency BCE loss receives misaligned supervision signals.
        Hungarian matching between predicted and ground-truth nodes is a
        promising fix currently under investigation.
\end{itemize}

\subsection{LLM Inference Latency and Deployment}

Phase~II with Qwen2.5-7B requires 2--10\,s per prediction on a 16\,GB GPU.
This is acceptable for offline batch processing but prohibitive for interactive
CAD tools where sub-second response is expected.  Three mitigation strategies
are available but not yet evaluated:

\begin{itemize}
  \item \textbf{Speculative decoding}: a small draft model generates candidate
        tokens that the main model verifies in parallel, reducing wall-clock
        time by 2--4$\times$ with no quality loss.
  \item \textbf{Smaller fine-tuned models}: Nemotron-Mini-4B
        (Section~\ref{sec:phase2_nemotron}) already reduces latency to
        $\sim$1\,s; fine-tuning it further with QLoRA is the next step.
  \item \textbf{Cached KV prefix}: the system prompt and few-shot examples are
        constant across queries and can be pre-computed and cached, eliminating
        redundant attention computation on the prompt portion.
\end{itemize}

\subsection{Evaluation Metric Completeness}

The metrics reported in Section~\ref{sec:metrics} measure coordinate-level
accuracy (MAE, RMSE, Hausdorff, Chamfer) and structural validity (PSR,
topology accuracy).  They do not directly measure \emph{engineering utility}:
whether the predicted midcurve, when used as a beam neutral axis in an FEA
solver, yields stress results within an acceptable tolerance of the full-solid
simulation.  Closing this loop—connecting neural midcurve prediction to
downstream simulation accuracy—is an important but currently unaddressed
validation dimension.

\section{Conclusion and Future Work}
\label{sec:conclusion}

\subsection{Summary}

This paper presented \textit{MidcurveNN}, a tri-paradigm neural framework for
computing the midcurve (medial axis) of 2D closed polygons.  By unifying three
fundamentally different representations—raster bitmaps, structured text, and
geometric graphs—under a single \emph{dimension-reduction} framing, the work
establishes a comprehensive reference point for learning-based approaches to
this long-standing geometric modelling problem.

\textbf{Phase~I} demonstrated that image encoder-decoder networks can learn the
profile-to-midcurve mapping from as few as 896 augmented image pairs.  Among
the four in-scope variants, the Convolutional Encoder-Decoder with CoordConv
and stride-2 skip connections produces the sharpest midcurve bitmaps with the
fewest stray pixels.  The Denoising Autoencoder's two-stage training provides
an effective and lightweight noise-suppression strategy.  The fundamental
lesson of Phase~I is that rasterisation provides a fast, low-barrier entry
point but its coordinate-precision ceiling is reached quickly; further
investment in the image paradigm yields diminishing returns.

\textbf{Phase~II} showed that large language models, when given an appropriate
geometric serialisation (BRep JSON), can reliably generate syntactically and
geometrically valid midcurve representations.  QLoRA fine-tuning of
Qwen2.5-7B on 993 BRep pairs achieves MAE\,=\,0.78, RMSE\,=\,1.24,
Hausdorff distance\,=\,2.1, and PSR\,=\,98\%—metrics that meet or exceed
the image baselines on every coordinate-level measure while also providing
explicit topology.  Nemotron-Mini-4B in a two-shot setting achieves 85.7\%
valid-JSON output and topology accuracy 0.83 on consumer hardware
($\approx$2\,GB VRAM), demonstrating that compact instruction-tuned models
are practical baselines for this task.  The BRep JSON format itself is a
transferable contribution: it is compact, human-readable, LLM-compatible, and
captures branched topology without special-purpose tokens.

\textbf{Phase~III} established the conceptually cleanest formulation: a
Graph Transformer that encodes the profile as a polygon graph and decodes the
midcurve as a skeleton graph, operating on exact floating-point coordinates
throughout.  The LaplacianPE\,$\to$\,TransformerConv\,$\to$\,TopKPooling\,%
$\to$\,decoder pipeline implements graph summarisation as a differentiable
end-to-end model, with a combined Chamfer\,+\,BCE loss that jointly supervises
coordinate accuracy and topological structure.  Comprehensive evaluation on a
larger benchmark is ongoing and will be reported separately.

\subsection{Graph-to-Graph Translation as the Next Frontier}

The central open problem that Phase~III begins to address—\emph{graph-to-graph
translation with different input and output topologies}—is, to the best of our
knowledge, underexplored in the geometric deep learning literature.  Standard
graph neural network benchmarks (molecular property prediction, social network
analysis, knowledge graphs) all involve fixed-topology tasks; none requires the
model to produce a structurally different graph.

We believe midcurve computation is an ideal \emph{canonical task} for
graph-to-graph learning: the input (closed polygon) and output (open/branched
polyline) are geometrically well-defined, the ground truth is unambiguous for
the shapes studied here, and the task difficulty scales controllably with shape
complexity.  Making this benchmark publicly available and attracting the
geometric deep learning community's attention to it is a near-term goal.

\subsection{Future Directions}

\paragraph{Diffusion-Based Graph Generation}
Denoising diffusion probabilistic models have achieved remarkable results in
image~\cite{ho2020ddpm} and molecule~\cite{hoogeboom2022edm} generation.
Adapting diffusion to the graph-to-graph midcurve problem—conditioning the
denoising process on the input profile graph—is a promising direction.
Diffusion naturally handles multi-modal output distributions (a profile with
near-symmetry may have two plausible midcurves) and produces well-calibrated
uncertainty estimates, which are valuable in a CAD context.

\paragraph{Integration of LLM Structural Reasoning with Geometric GNNs}
Phase~II and Phase~III represent complementary strengths: LLMs bring broad
structural reasoning and topology awareness from pre-training; Graph
Transformers bring geometric precision and coordinate-level accuracy.  A hybrid
architecture—using an LLM to predict the midcurve topology (number of
branches, junction degrees) and a Graph Transformer to predict exact vertex
coordinates conditioned on that topology—could combine both strengths.  The
BRep JSON intermediate representation studied in Phase~II is a natural
interface between the two modules.

\paragraph{Generalisation to Complex and Non-Standard Profiles}
Extending the dataset to include the \texttt{PhDdata/} complex profiles, concave
cross-sections (C, Z, Hat, U), and non-simply-connected shapes (concentric
rings) is essential for any claim of practical generalisability.  The augmented
BRep JSON pipeline (Phase~II) can generate these variants automatically from
manually authored base shapes; the bottleneck is authoring the base shapes
themselves, which requires domain knowledge of which midcurve is ``correct'' for
each profile family.

\paragraph{Downstream FEA Validation}
The ultimate test of a predicted midcurve is not a geometric distance metric
but engineering utility: does replacing the solid section with a beam element
on the predicted midcurve yield FEA stress results within an acceptable
tolerance of the full-solid simulation?  Coupling MidcurveNN predictions with
an open-source FEA solver (e.g.\ FEniCS or CalculiX) to produce end-to-end
simulation accuracy metrics would close the loop between neural geometry
research and practical engineering application.

\paragraph{Adaptive Pooling for Phase~III}
Resolving the fixed-$k$ TopKPooling limitation identified in
Section~\ref{sec:challenges} via shape-conditioned or attention-based adaptive
pooling is the highest-priority architectural improvement for Phase~III.
Together with Hungarian-matched adjacency supervision, this is expected to
unlock reliable quantitative evaluation and enable a fair head-to-head
comparison with Phase~II on the four-shape benchmark.

\subsection{Closing Remark}

Midcurve computation has been treated as a geometry-processing problem for
fifty years.  MidcurveNN reframes it as a \emph{learning} problem, and in
doing so, connects it to three of the most active areas of contemporary machine
learning: image generation, large language models, and geometric deep learning.
The three phases are not competing solutions but waypoints on a continuum from
approximate-but-accessible to exact-and-principled.  We hope this unified
framework encourages cross-pollination between the CAD/CAM community and the
deep learning community, and that the graph-to-graph formulation in particular
finds application beyond midcurves to the broader class of geometric
dimension-reduction problems.

%% ── Bibliography ────────────────────────────────────────────────────────────
% T_P12: Audit references; add MAT/Voronoi classics, UNet, Pix2Pix,
% QLoRA/PEFT, HuggingFace, Nemotron, Graph Transformer papers.
\bibliographystyle{IEEEtran}
\bibliography{references}

\end{document}
